\documentclass[a4paper,11pt]{article}
\usepackage[utf8]{inputenc}
\usepackage[T1]{fontenc}
\usepackage[french]{babel}
\usepackage[right=2.5cm, left=2.5cm, bottom=4cm, top=3cm]{geometry}
\usepackage[ddmmyyyy]{datetime}
\usepackage[table]{xcolor}
\usepackage{lmodern,mathptmx,changepage,titlesec,hyperref,listings,lstautogobble,graphicx,array,longtable,multirow,lipsum,tikz,shorttoc,enumitem,float,verbatim}
\usetikzlibrary{arrows,automata}
\usetikzlibrary{positioning}

\renewcommand{\rmdefault}{\sfdefault} %Utilisation de la police sans-serif ("Computer Modern Sans") pour la police roman
\renewcommand{\ttdefault}{pcr} 	%Utilisation d'une police "CourrierNew" pour la police monospaced (pour faire un listing manuel)
\linespread{1.15}				%Interligne

%Utilisation de liens colorés en bleu et soulignés
\hypersetup{colorlinks=true, urlcolor=blue, urlbordercolor=blue, linkcolor=black, linkbordercolor=white}
\makeatletter \Hy@AtBeginDocument{\def\@pdfborder{0 0 1} \def\@pdfborderstyle{/S/U/W 1}}\makeatother

\titlespacing*{\section} {0cm}{7ex plus 1ex minus .2ex}{1.5ex plus .2ex}
\titlespacing*{\subsection} {0cm}{4.5ex plus 1ex minus .2ex}{1.5ex plus .2ex}
\titleformat*{\section}{\LARGE\bfseries}
\titleformat*{\subsection}{\large\bfseries}
\titleformat*{\subsubsection}{\normalsize\bfseries}

\definecolor{darkgreen}{rgb}{0,0.8,0}
\definecolor{mygray}{rgb}{0.93,0.93,0.93}
\definecolor{mymauve}{rgb}{0.58,0,0.82}
\lstset{	
	language=Python,
	basicstyle=\small\ttfamily,
	backgroundcolor=\color{mygray},
	breaklines=true,
	breakatwhitespace=true,
	tabsize=3,
	frame=none,
	rulecolor=\color{black},
	keywordstyle=\color{blue}\bfseries,
	stringstyle=\color{orange},
	showstringspaces=false,
	commentstyle=\footnotesize\color{darkgreen},
	keepspaces=true,
	extendedchars=true,
	numbers=left,
	numberstyle=\tiny\color{lightgray},
	stepnumber=1,
	escapeinside={(@}{@)},
	autogobble=true,
	literate=
		{á}{{\'a}}1 {é}{{\'e}}1 {í}{{}}1 {ó}{{\'o}}1 {ú}{{\'u}}1
		{Á}{{\'A}}1 {É}{{\'E}}1 {Í}{{\'I}}1 {Ó}{{\'O}}1 {Ú}{{\'U}}1
		{à}{{\`a}}1 {è}{{\`e}}1 {ì}{{\`i}}1 {ò}{{\`o}}1 {ù}{{\`u}}1
		{À}{{\`A}}1 {È}{{\'E}}1 {Ì}{{\`I}}1 {Ò}{{\`O}}1 {Ù}{{\`U}}1
		{ä}{{\"a}}1 {ë}{{\"e}}1 {ï}{{\"i}}1 {ö}{{\"o}}1 {ü}{{\"u}}1
		{Ä}{{\"A}}1 {Ë}{{\"E}}1 {Ï}{{\"I}}1 {Ö}{{\"O}}1 {Ü}{{\"U}}1
		{â}{{\^a}}1 {ê}{{\^e}}1 {î}{{\^i}}1 {ô}{{\^o}}1 {û}{{\^u}}1
		{Â}{{\^A}}1 {Ê}{{\^E}}1 {Î}{{\^I}}1 {Ô}{{\^O}}1 {Û}{{\^U}}1
		{œ}{{\oe}}1 {Œ}{{\OE}}1 {æ}{{\ae}}1 {Æ}{{\AE}}1 {ß}{{\ss}}1
		{ç}{{\c c}}1 {Ç}{{\c C}}1 {ø}{{\o}}1 {å}{{\r a}}1 {Å}{{\r A}}1
		{€}{{e}}1 {£}{{\pounds}}1 {«}{{\guillemotleft}}1
		{»}{{\guillemotright}}1 {ñ}{{\~n}}1 {Ñ}{{\~N}}1 {¿}{{?`}}1
}

%Redéfinition de la taille de \Huge pour le titre du document
\makeatletter\renewcommand\Huge{\@setfontsize\Huge{37pt}{40}}\makeatother
\date{\today}


\title{\vspace{\fill}\textbf{\Huge Rapport}}
\author{
	Sonny Klotz - Idir Hamad - Younes Benyamna - Malek Zemni
	\vspace{2em}\\
	\textit{Projet M1 Informatique}\\\textit{Primalité}
	\vspace{2em}
}

\begin{document}
\pagenumbering{gobble}\clearpage
\maketitle\vspace{9em}
\begin{center}\includegraphics[scale=0.7]{logo.png}\end{center}
\begin{flushright}Module \textit{TER}\end{flushright}
\newpage
\tableofcontents
\newpage\clearpage\pagenumbering{arabic}

	\section*{Introduction}
	
	\paragraph{}Ce document est le compte-rendu final de notre travail qui s'inscrit dans le cadre du module \textit{Projet} du master informatique de l'\textit{\textbf{UVSQ}}.
	
	
	\section{Architecture de l'application}
		\subsection{Organigramme et données échangées}
		Cet organigramme préalablement établi, représente la décomposition en modules de l'application, ainsi que les informations qui circulent entre ces modules.
		
		\subsection{Fonctionnalités des modules}
		
	\section{Outils et langages de programmation}
		\subsection{Contraintes}
		Lors de la réalisation de l'application, des contraintes ont dû être respectées.
		Ces contraintes se sont avérées déterminantes pour les choix techniques que l'on a fait, et en l'occurrence, le choix des langages de programmation. 
		
		\subsection{Choix des outils et des langages}
		Les contraintes imposées et les exigences définies nous ont permis de fixer nos choix sur plusieurs langages de programmation qui vont interagir ensemble :
	
	\section{Fonctionnement de l'application}
	
	\section{Cryptosystèmes}
	
	\section{Tests de primalité}

		
	\section{Bilan technique du projet}
		Notre produit final, c'est à dire l'application, se comporte comme prévu : l'application est fonctionnelle, la liaison entre ses différents modules réussit bien et les différentes fonctionnalités fournissent le résultat attendu.		

		\subsection{Problèmes rencontrés}
			\subsubsection*{Problèmes résolus :} 
			Lors de la réalisation de l'application, on a été confrontés à plusieurs problèmes et points délicats, principalement des verrous techniques, qui ont perturbé le bon déroulement de notre travail :
				
			\subsubsection*{Problèmes non résolus :}
				Certains problèmes rencontrés n'ont pas été entièrement résolus. Ces problèmes ne sont pas déterminants pour l'acceptabilité de notre produit.
				
	\section{Organisation interne du groupe}
	Assignation des modules pour chaque membre du groupe :
	Cette répartition a été parfaitement respectée. Elle nous a permis de travailler efficacement et assez indépendamment, ce qui prouve que l'assignation des modules a été judicieusement faite. Nous sommes également restés en contact pendant toute la phase de développement pour s'entraider pour la prise en main des nouveaux outils.
	
	\section{Coûts}
	Ce tableau indique les coûts estimés et les coûts finaux, en nombre de lignes de code et pour chaque module :
	
	\section*{Conclusion}
		
\end{document}
